\sectionhead{\textbf{About this copy}}

This is a corrected copy of the paper cited above. The text is retyped from the printed paper.
It keeps the paper's wording and the numbering of its results, equations and footnotes, and a
marker such as \origpagemark{$N$} shows where page $N$ of the printed paper begins.

The corrections come from a formalization of the Graph Minors series in the Lean proof
assistant, \href{https://github.com/mahdi-jfri/GraphMinor}{github.com/mahdi-jfri/GraphMinor},
in which the results of this paper are proved. Where the formalization found that a statement
or a step of a proof is wrong as printed, this copy corrects it in place. Where the paper leaves a proof or a step to the reader, this
copy adds it, also in blue. Text \del{struck
through in red} is removed and text \fix{in blue} is inserted. Each change has a note, marked
\textcolor{fixcolor}{\sffamily F1}, \textcolor{fixcolor}{\sffamily F2}, \dots, that gives the
reason. The file
\href{https://github.com/mahdi-jfri/GraphMinor/blob/main/corrections/GM18.md}{corrections/GM18.md}
of the formalization describes each correction in full. Each note, and each item in the list below, begins with the kind of the change. A \kindname{misprint} is a wrong symbol, index or word. A \kindname{statement} is a result or a claim of the text that is false as printed. A \kindname{degenerate} is a statement that fails only in a trivial case, such as an empty graph or a parameter equal to~$0$. A \kindname{gap} is a step of a proof that does not follow as written. A \kindname{missing} is a proof that the paper omits, leaves to the reader or only sketches, and that this copy adds.

The corrections and added proofs, in the order of the pages:
\begin{enumerate}\setlength{\itemsep}{1pt}
\item \kindtag{fix:2.1-proof-added}\hyperref[fix:2.1-proof-added]{Proof of 2.1, page 79}: the proof is added.
\item \kindtag{fix:2.4-proof-added}\hyperref[fix:2.4-proof-added]{Proof of 2.4, page 80}: the proof is added.
\item \kindtag{fix:3.4-claim1}\hyperref[fix:3.4-claim1]{Claim (1) in the proof of 3.4, page 84}: $(A_k, B_{jk})$
becomes $(A_j, B_{jk})$.
\item \kindtag{fix:4.3-claim1}\hyperref[fix:4.3-claim1]{Claim (1) in the proof of 4.3, page 87}: a backslash in the
subscript of $\mathscr{L}$ becomes $i$, and $\subseteq$ becomes $\not\subseteq$.
\item \kindtag{fix:5.4-claim1}\hyperref[fix:5.4-claim1]{Claim (1) in the proof of 5.4, page 94}: $\Delta_j(E(A_j))$
becomes $\Delta_j(e(A_j))$.
\item \kindtag{fix:5.5-claim3}\hyperref[fix:5.5-claim3]{Claim (3) in the proof of 5.5, page 96}: $e(a)$ becomes $e(A)$.
\item \kindtag{fix:6.5-claim3-f0}\hyperref[fix:6.5-claim3-f0]{Claim (3) in the proof of 6.5, page 104}: $\tau \times (T, f)$
becomes $\tau \times (T, f_0)$.
\item \kindtag{fix:6.5-claim3-proof-added}\hyperref[fix:6.5-claim3-proof-added]{Claim (3) in the proof of 6.5, page 104}: the proof is
added.
\item \kindtag{fix:6.6-proof}\hyperref[fix:6.6-proof]{Proof of 6.6, page 106}: the bound $k$ on the number of roots is
taken over $\mathscr{C}^*$, not over $\mathscr{C}$.
\end{enumerate}
